\subsection{Exactly one T3-like role}

\begin{proposition}[Direct T3-like rescuer witness]
\label{prop:new-one-t3-terminal}
Assume $N_+=1$, exactly one V triangle is T3-like, no V triangle is
Vd1/Vd2, and at least one V-gap is present.  Then the skeleton is not covered.
\end{proposition}

\begin{proof}
Normalize the center midpoint to $M_0$.  Midpoint locality forces, after
reflection,
\[
 T_0\text{ T3-like},
 \qquad M_1\in T_0,
 \qquad T_1\text{ uniquely supercritical},
\]
while $T_2,T_3,T_4,T_5$ are nonsupercritical Vd0 roles.

Use the translated T3-like normal form at $V_0$.  Let its boundary reach on
$e_{5,0}$ be $a$, and write its positive interval on $r_1$, measured from
$V_1$ toward $O$, as
\[
 [c,u],
 \qquad c\le\frac12\le u.
\]
The O-side endpoint
\[
 P_T=(1-u)V_1
 \tag{6.41}
\]
is not in the open T3-like role.  The adjacent supercritical role contains
$V_1$ but misses $M_1$, so convexity excludes the O-side of $M_1$.  The four
other roles are Vd0.  Hence every V role misses $P_T$ and
\[
 \boxed{P_T\in U_C.}
 \tag{6.42}
\]

The normalized T3-like shape has parameters
\[
 1\le D\le\frac2{\sqrt3},
 \qquad E_0=\sqrt{4-3D^2},
 \qquad R_0=\frac{D+E_0}{2},
\]
and
\[
 c=\frac{D(1+a)-1}{R_0},
 \qquad
 u=1-\frac{R_0}{D}+a.
\]
Put
\[
 x=\frac{aD}{R_0},
 \qquad
 \theta=\frac{D-1}{R_0}.
\]
Then
\[
 c=x+\theta,
 \qquad
 0\le\theta\le2-\sqrt3,
\]
and the midpoint constraints are
\[
 \frac{1-4\theta+\theta^2}{2(1-2\theta)}
 \le x\le\frac12-\theta.
 \tag{6.43}
\]
Let $M_c^{\rm sup}$ be (6.8).  We claim
\[
 \boxed{
 a\le1-M_c^{\rm sup},
 \qquad
 \frac{a}{a+1-u}\le1-M_c^{\rm sup}.
 }
 \tag{6.44}
\]
For $0\le z\le1/2$, put
\[
 s(z)=\frac{z(2-z)}{1+z}.
\]
The defining equation for the supercritical envelope is
\[
 s(1-M_c^{\rm sup})=c.
\]
Since $s$ is increasing, it is enough to prove $s(x)\le x+\theta$.  After
multiplication by $1+x$, this becomes
\[
 Q_\theta(x):=2x^2+(\theta-1)x+\theta\ge0.
\]
For $0\le\theta\le1/5$, the feasible interval (6.43) lies to the right of
the vertex and substitution of its left endpoint gives
\[
 Q_\theta(x)\ge
 \frac{\theta(1-5\theta+11\theta^2-\theta^3)}
 {2(1-2\theta)^2}\ge0.
\]
For $1/5\le\theta\le2-\sqrt3$, the vertex lies in the feasible interval and
\[
 Q_\theta(x)\ge\frac{10\theta-1-\theta^2}{8}>0.
\]
Thus $x\le1-M_c^{\rm sup}$.  Since $D\ge R_0$, one has $a\le x$; moreover
\[
 \frac{a}{a+1-u}=\frac{aD}{R_0}=x.
\]
This proves (6.44).

Let $h_T$ be the boundary reach forced on $T_5$ from $V_5$ toward $V_0$.
If the companion center trace does not hide the endpoint $a$, then
\[
 h_T\ge1-a\ge M_c^{\rm sup}.
\]
In the only hiding configuration, the center near endpoint gives $k\le Wa$,
while (6.42) gives $\delta/R\ge1-u$.  Hence
\[
 R(a+1-u)\le a
\]
and the far center endpoint satisfies
\[
 R+\alpha\le\frac{a}{a+1-u}\le1-M_c^{\rm sup}.
\]
Thus in every case
\[
 \boxed{h_T\ge M_c^{\rm sup}.}
 \tag{6.45}
\]

The adjacent supercritical role has own-radial reach at least $c$, so
\[
 \boxed{B_1<M_c^{\rm sup}.}
 \tag{6.46}
\]
The four ordinary roles cover the center-free boundary path from $e_{1,2}$
to $e_{5,0}$.  Adding the edge-cover inequalities gives
\[
 \sum_{i=2}^5(A_i+B_i)
 \ge4+h_T-B_1>4,
\]
contrary to nonsupercriticality of all four roles.  For two gaps, the shorter
CE2 radial contradiction of Theorem~\ref{thm:new-ce2-short-ray} applies.
\end{proof}

\subsection{Exactly one Vd1/Vd2 role in CE2}

\begin{proposition}[Direct one-Vd placement assembly]
\label{prop:new-one-vd-assembly}
Assume CE2, $N_+=1$, exactly one Vd1/Vd2 role, and at least one V-gap.  Then
the skeleton is not covered.
\end{proposition}

\begin{proof}
Let $T_\sigma$ be the supercritical role and $T_\tau$ the Vd1/Vd2 role.
They are distinct.  If a further role has positive adjacent support, then
$N_++N_{\rm sp}\ge3$ and the skeleton-length theorem gives a contradiction.
Assume every other role is Vd0.

\smallskip
\noindent\emph{Case 1: $\sigma=0$ and $\tau$ is adjacent.}
Reflect to $\tau=1$.  Let $\rho_R,\rho_L$ be the residual boundary reaches
forced on the Vd role and on the far Vd0 path after removing the supercritical
trace and the two center intervals.  The Vd half-unit cap gives
\[
 \rho_R+\rho_L<\frac12.
 \tag{6.47}
\]
The signed CE2 inequalities and the two possible residual formulas give
\[
 \boxed{4\delta<\rho_L.}
 \tag{6.48}
\]
If the Vd role has support on $r_2$, its far endpoint satisfies
\[
 u_{1\to2}<1-\rho_L<1-\delta.
\]
For the ordinary role $T_2$, boundary coverage gives
\[
 A_2>\frac12+\rho_R,
 \qquad B_2\ge\rho_L.
\]
The quarter radial envelope therefore gives
\[
 C_2\le1-\frac{\rho_L}{4}<1-\delta.
\]
Thus neither local V trace reaches the center interval on $r_2$, whose
vertex-side entry is $1-\delta$.  All other roles are excluded by Vd0 and
diameter locality, so $r_2$ is uncovered.

\smallskip
\noindent\emph{Case 2: $\sigma=0$ and $\tau\in\{2,3,4\}$.}
Put $T=\alpha+\delta$.  Let $\rho_R,\rho_L$ be the boundary tails forced on
$T_\tau$.  Diameter transfer from the two center near endpoints gives
\[
 T<\min\{\rho_R,\rho_L\}.
 \tag{6.49}
\]
Every relevant center exit is at most $T$, whereas the Vd own-radial margin
gives
\[
 C_\tau<1-\min\{\rho_R,\rho_L\}.
\]
Hence the point
\[
 D_\tau=\min\{\rho_R,\rho_L\}V_\tau
\]
lies beyond both the Vd trace and the center trace.  Neighboring Vd0 and
nonlocal roles miss it, so it is an uncovered point.

\smallskip
\noindent\emph{Case 3: $\tau=0$.}
Midpoint rescue forces $\sigma\in\{1,5\}$; reflect to $\sigma=1$.
If $T_0$ is Vd2, its perimeter contribution is strictly below $1/3$.
Together with the center cap $<1/2$, the supercritical cap $2/\sqrt3$, and
four unit caps, this gives
\[
 \frac12+\frac13+\frac2{\sqrt3}+4<6,
\]
a Strategy~1 contradiction.

Assume $T_0$ is Vd1.  Write its supported interval on $r_1$ as
$[c,u]$, with $c\le1/2\le u$, and put $P=(1-u)V_1$.  This point is missed by
all V roles, hence lies in $U_C$.  The exact Vd1 corner form has parameter
$t\ge1$ and gives
\[
 c=\frac{t(1-b)}{t+1},
 \qquad
 u=\frac{\sqrt{t^2+t+1}-a-tb-1}{t}.
\]
Direct differentiation and the quadratic
\[
20c^2+(18\sqrt3-52)c+47-24\sqrt3
\]
prove
\[
 a\le1-M_c^{\rm sup},
 \qquad
 \frac{a}{a+1-u}\le1-M_c^{\rm sup}.
 \tag{6.50}
\]
Exactly as in the T3-like case, the center cannot hide the far boundary tail;
thus the terminal reach on $T_5$ is at least $M_c^{\rm sup}$, while the
adjacent supercritical role has $B_1<M_c^{\rm sup}$.  The four ordinary path
roles would have boundary-sum total greater than four, a contradiction.

\smallskip
\noindent\emph{Case 4: $\sigma\ne0$ and $\tau\ne0$.}
Midpoint rescue makes the two roles adjacent.  If the Vd role is Vd2, use the
preceding perimeter terminal.  If it is Vd1, the corrected two-chart
replacement produces six nonsupercritical Vd0 roles while preserving the
covered skeleton.  Recompute the output gap rank.  Rank zero is excluded by
the boundary-complete length theorem; rank one or two is excluded by
Proposition~\ref{prop:new-nplus-zero-gap-closures}.  No preservation of the
input gap rank is used.

The four cases are exhaustive.
\end{proof}
